\ticket{18}{Машинное обучение: обучение с учителем. Задача классификации, способы оценки качества решений. Классификация текстов: применяемые при обучении признаки текстов, методы обучения классификаторов (на примере двух методов)}

\begin{examplan}
    \item Дать определение обучения с учителем.
    \item Описать задачу классификации.
    \item Перечислить способы оценки качества классификации.
    \item Объяснить признаки текстов и два метода классификации текстов.
\end{examplan}

\subsection*{Короткая версия для письменного ответа}

\textbf{Обучение с учителем} --- машинное обучение на размеченной выборке пар $(x_i,y_i)$, где $x_i$ --- объект или его признаки, а $y_i$ --- правильный ответ. Модель учится по обучающим данным предсказывать ответы для новых объектов.

Основные задачи: классификация, где ответ --- метка класса, и регрессия, где ответ --- число.

\subsection*{Задача классификации}

\textbf{Классификация} --- построение алгоритма $a:X\to Y$, который относит объект к одному из заранее заданных классов $Y=\{c_1,\ldots,c_K\}$.

Типы:
\begin{itemize}
    \item бинарная классификация: два класса;
    \item многоклассовая классификация: более двух классов;
    \item многометочная классификация: объект может иметь несколько меток.
\end{itemize}

Процесс обучения: сбор данных, выделение признаков, обучение модели, проверка на тестовой выборке, применение к новым объектам.

\subsection*{Оценка качества}

Для бинарной классификации используют:
\begin{itemize}
    \item $TP$ --- истинно положительные;
    \item $TN$ --- истинно отрицательные;
    \item $FP$ --- ложноположительные;
    \item $FN$ --- ложноотрицательные.
\end{itemize}

\[
Accuracy=\frac{TP+TN}{TP+TN+FP+FN}.
\]

\[
Precision=\frac{TP}{TP+FP},\qquad Recall=\frac{TP}{TP+FN}.
\]

\[
F1=2\cdot\frac{Precision\cdot Recall}{Precision+Recall}.
\]

Accuracy плохо отражает качество при дисбалансе классов. В многоклассовых задачах метрики усредняют по классам: macro-average дает равный вес каждому классу, micro-average считает суммарные $TP$, $FP$, $FN$.

\subsection*{Классификация текстов}

Классификация текстов --- автоматическое отнесение документа к категории: спам/не спам, тема новости, тональность отзыва.

Текст нужно преобразовать в числовые признаки:
\begin{itemize}
    \item \textbf{Bag of Words}: учитывается наличие или частота слов, порядок игнорируется;
    \item \textbf{TF-IDF}: повышает вес слов, частых в документе, но редких в коллекции;
    \item \textbf{n-граммы}: учитывают локальный контекст;
    \item \textbf{лемматизация и стемминг}: уменьшают число разных форм слова;
    \item \textbf{стоп-слова}: часто удаляются как малоинформативные;
    \item \textbf{эмбеддинги}: плотные векторы слов или документов.
\end{itemize}

\subsection*{Наивный байесовский классификатор}

Идея: выбрать класс, максимизирующий апостериорную вероятность:
\[
a(x)=\arg\max_{y\in Y}P(y)P(x\mid y).
\]

В текстах используется наивное предположение: признаки-слова условно независимы при известном классе:
\[
P(x\mid y)\approx\prod_j P(x_j\mid y).
\]

Для мультиномиальной модели вероятности слов оцениваются по частотам в документах класса. Обязательно используют сглаживание, чтобы неизвестные слова не давали нулевую вероятность.

Плюсы: простота, скорость, хорошая базовая модель для текстов. Минусы: предположение независимости слов грубое, взаимодействия между словами не учитываются.

\subsection*{Логистическая регрессия}

\textbf{Логистическая регрессия} --- линейный вероятностный классификатор. Для бинарной задачи:
\[
P(y=1\mid x)=\sigma(w^Tx+b)=\frac{1}{1+e^{-(w^Tx+b)}}.
\]

Параметры обучаются минимизацией кросс-энтропии, часто с L1- или L2-регуляризацией. Для многоклассовой классификации используется softmax-регрессия.

Плюсы: дает вероятности, интерпретируема по весам признаков, хорошо работает с TF-IDF. Минусы: линейная граница решений, поэтому сложные нелинейные зависимости требуют дополнительных признаков или других моделей.

\begin{important}
    \item Обучение с учителем требует размеченных примеров.
    \item Классификация предсказывает дискретную метку класса.
    \item Precision и recall важны при дисбалансе и разной цене ошибок.
    \item Для текстов ключевой этап --- превращение текста в признаки.
    \item Наивный Байес быстр и прост; логистическая регрессия сильная линейная базовая модель.
\end{important}
